%%
%% This is file `sample-sigconf-authordraft.tex',
%% generated with the docstrip utility.
%%
%% The original source files were:
%%
%% samples.dtx  (with options: `all,proceedings,bibtex,authordraft')
%% 
%% IMPORTANT NOTICE:
%% 
%% For the copyright see the source file.
%% 
%% Any modified versions of this file must be renamed
%% with new filenames distinct from sample-sigconf-authordraft.tex.
%% 
%% For distribution of the original source see the terms
%% for copying and modification in the file samples.dtx.
%% 
%% This generated file may be distributed as long as the
%% original source files, as listed above, are part of the
%% same distribution. (The sources need not necessarily be
%% in the same archive or directory.)
%%
%%
%% Commands for TeXCount
%TC:macro \cite [option:text,text]
%TC:macro \citep [option:text,text]
%TC:macro \citet [option:text,text]
%TC:envir table 0 1
%TC:envir table* 0 1
%TC:envir tabular [ignore] word
%TC:envir displaymath 0 word
%TC:envir math 0 word
%TC:envir comment 0 0
%%
%% The first command in your LaTeX source must be the \documentclass
%% command.
%%
%% For submission and review of your manuscript please change the
%% command to \documentclass[manuscript, screen, review]{acmart}.
%%
%% When submitting camera ready or to TAPS, please change the command
%% to \documentclass[sigconf]{acmart} or whichever template is required
%% for your publication.
%%
%%
\documentclass[sigconf]{acmart}
%%
%% \BibTeX command to typeset BibTeX logo in the docs
\AtBeginDocument{%
  \providecommand\BibTeX{{%
    Bib\TeX}}}

%% Rights management information.  This information is sent to you
%% when you complete the rights form.  These commands have SAMPLE
%% values in them; it is your responsibility as an author to replace
%% the commands and values with those provided to you when you
%% complete the rights form.
\setcopyright{none}
\copyrightyear{2026}
\acmYear{2026}
\acmDOI{XXXXXXX.XXXXXXX}
%% These commands are for a PROCEEDINGS abstract or paper.
\acmConference[PREPRINT]{ICSE'26 Companion. Please note that this document is a preliminary version}{2026}{Accepted at ICSE'26 Track ``Software architecture BoF''}
%%
%%  Uncomment \acmBooktitle if the title of the proceedings is different
%%  from ``Proceedings of ...''!
%%
%%\acmBooktitle{Woodstock '18: ACM Symposium on Neural Gaze Detection,
%%  June 03--05, 2018, Woodstock, NY}
%\acmISBN{978-1-4503-XXXX-X/2018/06}
\acmISBN{XXX-X-XXXX-XXXX-X/XXXX/XX}


%%
%% Submission ID.
%% Use this when submitting an article to a sponsored event. You'll
%% receive a unique submission ID from the organizers
%% of the event, and this ID should be used as the parameter to this command.
%%\acmSubmissionID{123-A56-BU3}

%%
%% For managing citations, it is recommended to use bibliography
%% files in BibTeX format.
%%
%% You can then either use BibTeX with the ACM-Reference-Format style,
%% or BibLaTeX with the acmnumeric or acmauthoryear sytles, that include
%% support for advanced citation of software artefact from the
%% biblatex-software package, also separately available on CTAN.
%%
%% Look at the sample-*-biblatex.tex files for templates showcasing
%% the biblatex styles.
%%

%%
%% The majority of ACM publications use numbered citations and
%% references.  The command \citestyle{authoryear} switches to the
%% "author year" style.
%%
%% If you are preparing content for an event
%% sponsored by ACM SIGGRAPH, you must use the "author year" style of
%% citations and references.
%% Uncommenting
%% the next command will enable that style.
%%\citestyle{acmauthoryear}



%%%%%%%%%%%%%%%%%
% Additional packages and commands
\usepackage{todonotes}
\newcommand{\inlinetodo}[1]{\todo[inline]{#1}}

\theoremstyle{definition}
\newtheorem{definition}{Definition}

%%%%%%%%%%%%%%%%%



%%
%% end of the preamble, start of the body of the document source.
\begin{document}

%%
%% The "title" command has an optional parameter,
%% allowing the author to define a "short title" to be used in page headers.
%\title{Towards "Creative Systems": the coming age of AI-enabled systems for quality achievement at runtime}
\title{Toward architecting self-coding information systems}

%%
%% The "author" command and its associated commands are used to define
%% the authors and their affiliations.
%% Of note is the shared affiliation of the first two authors, and the
%% "authornote" and "authornotemark" commands
%% used to denote shared contribution to the research.

% \author{Rodrigo Falc\~{a}o}
% %\authornote{Both authors contributed equally to this research.}
% \email{rodrigo.falcao@iese.fraunhofer.de}
% \orcid{0000-0003-1222-0046}
% \author{Frank Elberzhager}
% %\authornotemark[1]
% \email{frank.elberzhager@iese.fraunhofer.de}
% \affiliation{%
%   \institution{Fraunhofer IESE}
%   \city{Kaiserslautern}
%  % \state{Ohio}
%   \country{Germany}
% }


\author{Rodrigo Falc\~{a}o}
\orcid{0000-0003-1222-0046}
\affiliation{%
  \institution{Fraunhofer IESE}
  \city{Kaiserslautern}
  \country{Germany}}
\email{rodrigo.falcao@iese.fraunhofer.de}

\author{Frank Elberzhager}
\orcid{0000-0002-8748-3927}
\affiliation{%
  \institution{Fraunhofer IESE}
  \city{Kaiserslautern}
 % \state{Ohio}
  \country{Germany}
}
\email{frank.elberzhager@iese.fraunhofer.de}

\author{Karthik Vaidhyanathan}
\affiliation{%
  \institution{IIIT Hyderabad}
  \city{Hyderabad}
  \country{India}}
\email{karthik.vaidhyanathan@iiit.ac.in}


%%
%% By default, the full list of authors will be used in the page
%% headers. Often, this list is too long, and will overlap
%% other information printed in the page headers. This command allows
%% the author to define a more concise list
%% of authors' names for this purpose.
\renewcommand{\shortauthors}{Falc\~{a}o et al.}

%%
%% The abstract is a short summary of the work to be presented in the
%% article.
\begin{abstract}


In this extended abstract, we propose a novel research topic in the field of agentic AI, which we refer to as self-coding information systems. These systems will be able to dynamically adapt their structure or behavior by evaluating potential adaptation decisions, generate source code, test, and (re)deploy their source code autonomously, at runtime, reducing the time to market of new features. Here we motivate the topic, provide a formal definition of self-coding information systems, discuss some expected impacts of the new technology, and indicate potential research directions.



  
\end{abstract}

%%
%% The code below is generated by the tool at http://dl.acm.org/ccs.cfm.
%% Please copy and paste the code instead of the example below.
%%
\begin{CCSXML}
<ccs2012>
   <concept>
       <concept_id>10011007.10011074.10011092.10011782</concept_id>
       <concept_desc>Software and its engineering~Automatic programming</concept_desc>
       <concept_significance>500</concept_significance>
       </concept>
   <concept>
       <concept_id>10010147.10010178</concept_id>
       <concept_desc>Computing methodologies~Artificial intelligence</concept_desc>
       <concept_significance>500</concept_significance>
       </concept>
   <concept>
       <concept_id>10011007.10011074.10011075.10011077</concept_id>
       <concept_desc>Software and its engineering~Software design engineering</concept_desc>
       <concept_significance>500</concept_significance>
       </concept>
   <concept>
       <concept_id>10011007.10011074.10011075.10011078</concept_id>
       <concept_desc>Software and its engineering~Software design tradeoffs</concept_desc>
       <concept_significance>500</concept_significance>
       </concept>
 </ccs2012>
\end{CCSXML}

\ccsdesc[500]{Software and its engineering~Automatic programming}
\ccsdesc[500]{Computing methodologies~Artificial intelligence}
\ccsdesc[500]{Software and its engineering~Software design engineering}
\ccsdesc[500]{Software and its engineering~Software design tradeoffs}

%%
%% Keywords. The author(s) should pick words that accurately describe
%% the work being presented. Separate the keywords with commas.
\keywords{agentic AI, generative AI, software architecture, LLM, runtime}
%% A "teaser" image appears between the author and affiliation
%% information and the body of the document, and typically spans the
%% page.


\received{9 January 2026}
\received[revised]{XX January 2026}
\received[accepted]{XX January 2026}

%%
%% This command processes the author and affiliation and title
%% information and builds the first part of the formatted document.
\maketitle

\section{Motivation}\label{sec:motivation}

%Generative AI, in particular large language models (LLMs), has become an absolutely pervasive topic in software engineering. When it comes to software architecture, 


We have experienced unprecedented interest in artificial intelligence (AI) in the current decade. This can be asserted in the software engineering community, in general, and the software architecture community, specifically. Last year's ICSE has featured, once again, AI in multiple tracks, sessions, workshops, and keynotes\footnote{\url{https://conf.researchr.org/program/icse-2025/program-icse-2025/}}. The same can be observed in the software architecture community: ICSA 2025, held in Odense, had the theme ``Architecting for the next generation of intelligent systems''. The main conference featured AI in two keynotes and multiple sessions.

Generative AI, especially language models, plays an important role in the current ``AI spring''. These models have shown impressive capabilities in generating human-like content for virtually any task in any domain. Furthermore, base models can be (and have been) fine-tuned to excel on specific tasks such as writing code in a given programming language, automating configuration, text translation, text summarization, engaging in conversation, performing sentiment analysis, and doing image classification and generation, to name a few. This has allowed exploration of various use cases across domains. %Among the most popular LLM use cases is implementing specialized AI assistants (chatbots).

When applied to software architecture, generative AI can be viewed as a two-fold field. One is concerned about how generative AI can support software architects to perform architecture activities at development time (e.g., quality assessment of UML diagrams \cite{de2025llm}, generation of architectural design decisions \cite{dhar2024can}, documentation evaluation, code assistance, code generation, and test case creation  \cite{jahic2024state}).
On the other hand, there has been a growing interest in understanding how to architect systems that implement generative-AI-based components to achieve certain goals at runtime -- a comparatively less explored aspect of generative AI in software architectures~\cite{esposito2025generative}. We can divide this runtime aspect into two parts: using genAI to improve \textit{functional aspects} of the systems and using genAI components to improve \textit{quality aspects} of the systems. The former, on the one hand, has gained more attention, especially through the introduction of the ``retrieval augmented generation'' (RAG) in 2020 \cite{lewis2020retrieval}. From a purely conceptual perspective, RAG realized and improved a long-ago idea for ``intelligent information systems'' \cite{lebowitz1983intelligent}. Powered by LLMs, the concept thrived, and research on RAG has been frequently found, from reference architectures (e.g., \cite{xu2025design}) to practical use cases (e.g., \cite{adnan2025leveraging}).

On the other hand, research on genAI at runtime to address quality aspects is now taking its first steps. For example, a recent experiment with LLMs showed that they can be very effective in addressing interoperability \cite{falcao2025evaluating}. One of the strategies used was code generation \textit{at runtime} (which makes sense as LLM-based code has already been vastly adopted at design time to support architects and developers): in the experiment, the system used LLMs to generate code and, afterward, the code was cleaned, compiled, deployed, tested, and used to convert data from an unknown representation to a desired target representation -- all this on the fly. This experience contributes to a vision of what we regard as a coming research topic: the architecting of \textbf{self-coding information systems}.



% \inlinetodo{
% •	Status of development of [Gen]AI technology
% •	Its applications to and effectiveness in code-related tasks
% •	The increasing need for software-based autonomy
% •	An envisioned future: “creative software systems”* --> self-coding systems
% }


\section{Self-coding information systems}\label{sec:whatis}

% \begin{definition}
%     A self-coding information system is a software-based system that has the ability to implement functional requirements and/or address quality requirements by generating the source code of at least one of its subsystems and using the generated source code at runtime.
% \end{definition}
We define a self-coding information system as \textit{an information system whose software part can implement functional requirements or address quality requirements by generating the source code of at least one of its subsystems and using the generated source code at runtime}. At the core of the definition is the capability of these systems to generate (or regenerate) their source code. This idea builds on top of existing definitions that hint at autonomy, such as autonomous systems \cite{rasmus2023autonomous}, context-aware systems \cite{schilit1994context}, autonomic computing \cite{kephart2003vision} (in this case, the authors elaborate on a broad vision of self-* systems that includes autonomous modification of software, but in terms of applying software patches, not self-coding) and self-adaptive systems (see e.g., \cite{cheng2008software}, where the authors do mention adaptation that changes the source code but through dynamic source-code weaving, not actual code generation). While our definition does not prescribe how these systems achieve their self-coding capabilities, agents based on language models (especially those tuned for coding tasks) are nowadays the most straightforward option; therefore, we assume it in this paper.

Self-coding information systems show a high degree of modifiability. As defined by Bass et al. in \cite{bass2013software}, ``[m]odifiability is about change'' and must be planned by architects considering what can change, the likelihood of the change, when and by whom the change is made, and the change costs. In self-coding information systems, modifications can be done to the source code, at runtime (reduced time to market), by the system itself (autonomy), and cheaper as if it were done by humans. Conversely, in traditional software-based systems, configuration is a common practice for modifiability at runtime; however, it has a limited reach, for the more configurable a system is, the more complex its architecture becomes. Modifications that require change in the source code are traditionally performed at development time.

%therefore have great potential where these qualities are desired but may be difficult to achieve (for example, due to their costs). In Section~\ref{sec:motivation} we briefly discussed an example showing how systems could self-code to achieve interoperability. Interoperability is challenging in many domains not because of technical barriers but because of its implementation costs (e.g., \cite{stegemann2019interoperability}),   This relates to a broader quality  flexibility: For an example of the self-coding of a functional requirement, we can imagine a system that has information about their data model and can implement new reports autonomously, as requested by their users, 

%The first appearance of the term ``self-coding systems'' is found in a paper from 1956 \cite{grems1956truly}. Back then, the authors elaborated on the idea of raising the abstraction level of programming activities -- at a time when the ``physical characteristics of one special computer'', as they mention, should be considered. In their concrete case, they created an abstraction layer that allowed programmers to focus on writing algebraic expressions  (which are relevant in their application domain) and the integrated system (1) read the expressions and translated them into the machine language, and (2) computed the result from input data using the coded machine language.

%Apart from the aforementioned paper, the term ``self-coding system'' does not bring any result on the major computer-science-related scientific search engines. As of January 8, 2026, ACM Digital Library\footnote{\url{https://dl.acm.org/}} found only that paper, whereas neither IEEExplore\footnote{\url{https://ieeexplore.ieee.org/}} nor Scopus\footnote{\url{https://www.scopus.com/}} returns any result for the exact phrase. Google Scholar\footnote{\url{https://scholar.google.com/}} generated 34 results, most not related to computer science. The few computer science related documents refer to self-coding in the context of encoders and decoders in the machine learning field \cite{moraboena2020deep, fu2025unified}.

% \begin{definition}
% Self-coding systems as self-developing
% \end{definition}

% Autonomous systems


% \inlinetodo{
% •	Defining creative software systems
% [Gen]AI-enabled (nowadays, powered by language models)
% Implement functional requirements
% Address quality requirements
% Self-coding, self-testing, self-healing, self-monitoring, self-deploying…
% •	Contrasting with similar pre-existing definitions (e.g., autonomous systems, context-aware systems, creative systems, etc.)
% }


% \cite{truszkowski2006autonomous} is an example.

\section{Expected impact}


% \inlinetodo{
% •	Impact on the software engineering role(s)
% Software engineers who build and evolve systems vs. Software engineers who supervise and “grows” systems (ref. ICSA 2025 Keynote from Helena --> I think there is also a paper published at EUROMICRO 2025 from Helena and Jan Bosch that contains the content of the keynote)
% human developers as a luxus?
% loosing control / ownership over the code
% Where is the human in the loop (if any?)
% - new UI/UX concepts (chat-based interfaces?)
% engineerng a "meta system"
% •	Social impact
% }

%Based on the current development of how genAI is used by different roles in the software engineering cycle, it is not completely foreseeable how strong the impact will be. However, we already see that especially implementation and testing tasks are heavily supported by genAI up to tasks that are completely automated. Still, complete automation seems a bit farer away as the final reliability on the results genAI produces is often not 100 percentage. There is something like a shift-left in the software development pipeline, i.e. a stronger automation is in phases such as implementation, and earlier phases such as design or architecture become more attention and a stronger quality gate: When the basis for automated implementation with genAI is documents created earlier such as requirements or architecture designs, the quality has even be higher so that genAI does not create results based on bad quality documents. As human beings might consider their experience and knowledge to compensate for not ideal quality of the input, genAI cannot compensate the same way. A stronger consideration of context seems one way to compensate. If this or other  approaches are further researched, the automation level will probably grow, and roles might at least change from their profile, i.e. become more of a "manager of code" instead of writing all the code. Agents in this case become more like a pair, similar to what nowadays is done in techniques like pair programming with two human beings. Such AI components might, on the next level, organize themselves and evolve autonomously over time.

We focus on three aspects to classify the impact of the emergence and establishment of self-coding information systems: %into the following categories:

\paragraph{Architectural trade-offs} Certain architectural trade-offs for the improved modifiability might appear. For example, the analyzability of self-coding information systems, as self-generated code may be harder to understand; functional correctness can also come into question as language models are non-deterministic; resource utilization is increased by the energy consumption required for performing inferences; installability is also impacted as, depending on the size of the models and the system usage requirements, special hardware (e.g., GPUs) may be needed.

\paragraph{Software engineering roles} Existing roles must be adapted to this new class of systems. For example, requirements engineering and software architecture activities will expand to a metalevel: beyond specifying and implementing a certain desired function, software engineers will be concerned about specifying and implementing an autonomous component that can implement, test, and deploy a class of desired functions. Considering interoperability, for instance again, the LLM-based agent is not general purpose but tailored for the concrete and specific interoperability task. In some cases, user interfaces may include natural language interaction to trigger the self-coding capabilities (see, for instance, OpenAI's App SDK \cite{openai2025introducingApps}, which promises to allow users to ``talk'' with their apps).

\paragraph{Social and organizational aspects} Finally, social and organizational impacts may emerge from scenarios where self-coding information systems exist. On the one hand, companies might be able to evolve their systems faster. We assume that the level of autonomy will grow and consider more and more software system qualities, such as interoperability, maintainability, or testability. As the general idea is not restricted to a domain or sector, to a certain size of a software system, to a certain programming language, or something else, the general potential is rather high, and thus the impact. Productivity might grow significantly. On the other hand, would human developers become a luxury in small and medium-size organizations if self-coding information systems advanced? Would the role of software developer be (at least partially) replaced with an advanced user who could chat with the system to evolve its functionalities?
    

% \begin{itemize}
%     \item \textbf{Software engineering roles} must be adapted to this new class of systems. For example, requirements engineering and software architecture activities will expand to a metalevel: beyond specifying and implementing a certain desired function, software engineers will be concerned about specifying and implementing an autonomous component that can implement, test, and deploy a class of desired functions. Considering interoperability, for instance again, the LLM-based agent is not general purpose but tailored for the concrete and specific interoperability task. In some cases, user interfaces may include natural language interaction to trigger the self-coding capabilities (see, for instance, OpenAI's App SDK \cite{openai2025introducingApps}, which promises to allow users to ``talk'' with their apps).
%     \item Certain \textbf{architectural trade-offs} for the improved modifiability might appear. For example, the analyzability of self-coding information systems, as self-generated code may be harder to understand; functional correctness can also come into question as language models are non-deterministic; resource utilization is increased by the energy consumption required for performing inferences; installability is also impacted as, depending on the size of the models and the system usage requirements, special hardware (e.g., GPUs) may be needed.

%     \item Finally, \textbf{social and organizational impacts} may emerge from scenarios where self-coding information systems exist. One the one hand, companies might be able to evolve their systems faster. We assume that the level of autonomy will grow and consider more and more software system qualities, such as interoperability, maintainability, or testability. As the general idea is not restricted to a domain or sector, to a certain size of a software system, to a certain programming language, or something else, the general potential is rather high, and thus the impact. Productivity might grow significantly. On the other hand, would human developers become a luxury in small and medium-size organizations if self-coding information systems advanced? Would the role of software developer be (at least partially) replaced with an advanced user who could chat with the system to evolve its functionalities?
    
% \end{itemize}

%In the area of genAI at runtime, the step towards self coding might be faster. 

%Right now, genAI is able to create solutions directly (at development time), or implement a more deterministic solution. The same argument than before still holds, i.e. additional mechanisms for a better reliability are still under research. However, we assume that the level of autonomy will grow and consider more and more software system qualities, such as interoperability, maintainability, or testability. As the general idea is not restricted to a domain or sector, to a certain size of a software system, to a certain programming language, or something else, the general potential is rather high, and thus the impact. Productivity might grow significantly.

\section{Research directions for self-coding information systems}

% \inlinetodo{
% •	Challenges
%     - reliability
% •	Limitations
%     - doesnt work all the time
% •	Research directions
%     - reliability strategies
%     - architectural patterns
%     - self-architecting systems --> systems that are able to "explain themselves"?
%     - economic analysis of self-coding systems
% }

A big challenge in today's solution is the ``missing last step'', i.e., solutions with genAI often work, but the final reliability is often missing. This is usually not acceptable from a quality perspective, as it is unclear when the results are correct; respectively, this has to be found out with additional effort or mechanisms to compensate. Researching \textbf{reliability strategies} for self-coding information systems is an essential research direction.

%\paragraph{Architectural patterns}
\textbf{Reference architectures} must be designed (or abstracted from concrete architectures as patterns) describing common ways of realizing self-coding information systems. These patterns will be informative from different perspectives. For example, they can guide deployment decisions or help limit the scope of the self-coding functionalities.

%\paragraph{Maintainability}
Another aspect is how good self-coding information systems are considered from the perspective of \textbf{maintainability} in general. We face the situation that auto-generated code is sometimes difficult to understand, compromising the ability of a human developer to evolve or fix a bug in it. Furthermore, self-coding information systems are not exclusively self-coded, i.e., human developers will also change the code. Therefore, proper workflows to support the coexistence with human developers are needed. Note that this goes beyond current practices involving the usage of coding assistants, as they are not part of the target system and do not update the target system at runtime.

As self-coding information systems require special hardware, and also considering the inference costs, it is necessary to develop simple ways of determining whether a self-coding capability is adequate or is inadequate for a certain aspect of a given system from an \textbf{economic perspective}. Returning to the interoperability example of Section~\ref{sec:motivation}: given that a self-coding component and a team of developers would both produce the same functional correctness outcome, which strategy would be more economical? Several variables can play a role in such an equation.

%\paragraph{Self-architecting}
Finally, self-coding information systems could evolve to become not only self-coding but also \textbf{self-architecting systems}, which, for example, document their design decisions and at least partially could also review their architectures.



%Furthermore, it is still unclear what is the ideal context that we should provide genAI in order to fulfill the tasks at its best. 



%%
%% The acknowledgments section is defined using the "acks" environment
%% (and NOT an unnumbered section). This ensures the proper
%% identification of the section in the article metadata, and the
%% consistent spelling of the heading.
\begin{acks}
To our colleagues Rafael Capilla, Pablo Antonino, and Ivan Compagnucci, whose ideas in the co-organization of the SAGAI workshop have contributed to this manuscript (R.F. and F.E.). AI tools were used for proofreading.
\end{acks}

%%
%% The next two lines define the bibliography style to be used, and
%% the bibliography file.
\bibliographystyle{ACM-Reference-Format}
\bibliography{refs}


%%
%% If your work has an appendix, this is the place to put it.


\end{document}
\endinput
%%
%% End of file `sample-sigconf-authordraft.tex'.
